\chapter{Спин-накрытие и мост к сфере проекторного дефекта}

Предыдущая глава локализовала единственный недостающий топологический
мост: необходимо связать проекторное поле со значениями в
$\mathbb{RP}^2$ с канонической сферой направлений
$\operatorname{Spin}(3)/\operatorname{Spin}(2)$. В проекте уже встречалось
утверждение о двух подъёмах $n$ и $-n$, но оно применялось после выбора
ежа. Теперь отдельно проверяется, что именно следует из топологии и
вращательной эквивариантности, а что по-прежнему является условием сектора.

\section{Две сферы и одно проекторное отображение}

Пусть окружающая предполагаемый точечный дефект пространственная сфера
записана как
\begin{equation}
 S^2_{\rm sp}
 \simeq SO(3)/SO(2)
 \simeq \operatorname{Spin}(3)/\operatorname{Spin}(2).
 \label{eq:v5-spin-bridge-spatial-sphere}
\end{equation}
Пространство неориентированных ранг-один осей равно
\begin{equation}
 \mathbb{RP}^2_{\rm fam}
 \simeq SO(3)/O(2),
 \label{eq:v5-spin-bridge-projective-target}
\end{equation}
где $O(2)$ понимается как стабилизатор неориентированной оси внутри
$SO(3)$.

Каноническое отображение имеет вид
\begin{equation}
 q:S^2_{\rm sp}\longrightarrow\mathbb{RP}^2_{\rm fam},
 \qquad
 q(\mathbf n)=[\mathbf n],
 \qquad
 P(\mathbf n)=\mathbf n\mathbf n^{\mathsf T}.
 \label{eq:v5-spin-bridge-quotient-map}
\end{equation}
Это ровно проекторный ёж, уже вычисленный в главе о точечном дефекте.

\section{Единственность вращательно-эквивариантного ежа}

Рассмотрим непрерывное $SO(3)$-эквивариантное отображение
\begin{equation}
 f:S^2\longrightarrow\mathbb{RP}^2.
\end{equation}
Оно полностью определяется значением $f(\mathbf e_3)$. Стабилизатор
$\mathbf e_3$ есть группа вращений $SO(2)$ вокруг третьей оси, поэтому
$f(\mathbf e_3)$ должно быть осью, неподвижной относительно всех этих
вращений. Единственная такая точка в $\mathbb{RP}^2$ есть
$[\mathbf e_3]$. Следовательно,
\begin{equation}
 f(R\mathbf e_3)=R f(\mathbf e_3)=[R\mathbf e_3],
\end{equation}
то есть $f=q$.

Тот же вывод следует из однородных пространств: включение стабилизаторов
\begin{equation}
 SO(2)\subset O(2)
\end{equation}
задаёт естественную проекцию
$SO(3)/SO(2)\to SO(3)/O(2)$. Морфизмы однородных $G$-пространств и
орбитальная категория систематически используются в эквивариантной
гомотопической теории; см. \path{arXiv:1903.03152} и
\path{arXiv:2012.11912}.

\begin{theorem}[Единственность эквивариантного проекторного профиля]
Каноническое отображение $q(\mathbf n)=[\mathbf n]$ является единственным
$SO(3)$-эквивариантным отображением
$S^2\to\mathbb{RP}^2$. В частности, постоянная карта нулевого сектора не
является $SO(3)$-эквивариантной, поскольку в $\mathbb{RP}^2$ нет точки,
фиксированной всей группой $SO(3)$.
\end{theorem}

Это утверждение не требует выбора локального координатного разреза и не
вводит нового семейного $SU(2)$-дублета. Оно использует только уже заданное
действие пространственных вращений на векторе $\mathbf n$ и его
проекторном чтении.

\section{Подъём накрытия и знак заряда}

Проекция
\begin{equation}
 \pi:S^2\longrightarrow\mathbb{RP}^2,
 \qquad \pi(\mathbf n)=[\mathbf n],
\end{equation}
является универсальным двулистным накрытием. Поскольку окружающая сфера
$S^2_{\rm sp}$ односвязна, любое отображение
$S^2_{\rm sp}\to\mathbb{RP}^2$ допускает глобальный подъём после выбора
одной точки. Два подъёма отличаются антиподным преобразованием.

Для канонического $q$ эти подъёмы равны
\begin{equation}
 \widetilde q_+(\mathbf n)=\mathbf n,
 \qquad
 \widetilde q_-(\mathbf n)=-\mathbf n.
 \label{eq:v5-spin-bridge-two-lifts}
\end{equation}
Их степени равны
\begin{equation}
 \deg\widetilde q_+=+1,
 \qquad
 \deg\widetilde q_-=-1.
 \label{eq:v5-spin-bridge-lift-degrees}
\end{equation}
Именно так задаётся заряд ежа в нематическом порядке с целевым
$\mathbb{RP}^2$; два подъёма дают внешний и внутренний ежи противоположных
знаков. См. \path{arXiv:1107.1169}, \path{arXiv:2203.07311} и обзор
\path{arXiv:1706.06861}.

\section{Индуцированная хопфова линия}

На ориентированной целевой сфере рассмотрим положительную собственную
линию спинорного проектора
\begin{equation}
 P_+(\mathbf n)=\frac12
 \left(I+\mathbf n\cdot\boldsymbol\sigma\right).
\end{equation}
Это хопфова линия $H\to S^2$ с $c_1(H)=1$. Её обратные образы вдоль двух
подъёмов имеют
\begin{equation}
 L_+=\widetilde q_+^*H,
 \qquad c_1(L_+)=+1,
\end{equation}
\begin{equation}
 L_-=\widetilde q_-^*H,
 \qquad c_1(L_-)=-1,
 \qquad L_-\simeq L_+^*.
 \label{eq:v5-spin-bridge-hopf-pair}
\end{equation}
Главные $U(1)$-носители этих линий имеют полное пространство $S^3$.

Таким образом, прежнее назначение
\begin{equation}
 E\longmapsto L_+,
 \qquad
 E^*\longmapsto L_-
\end{equation}
теперь получает точное основание внутри вращательно-эквивариантного
проекторного сектора. Моритовская градуировка не создаёт величину класса
Черна; величина $|c_1|=1$ приходит от степени единственного
эквивариантного ежа, а градуировка согласует его два знака с $E/E^*$.

\begin{theorem}[Замыкание spin-cover bridge в секторе ежа]
Для канонического $SO(3)$-эквивариантного проекторного дефекта карта
$S^2_{\rm sp}\to\mathbb{RP}^2_{\rm fam}$ имеет два глобальных подъёма
степеней $+1$ и $-1$. Их спинорные собственные линии образуют хопфову пару
$L/L^*$ с классами Черна $+1/-1$. Следовательно, коэффициентный проектор
ранга пятнадцать даёт классы $+15/-15$ без дополнительного выбора линии.
\end{theorem}

\section{Почему это ещё не полное принуждение материи}

Полученный мост устраняет одну реальную недосказанность, но не разрешает
подменять условие сектора выводом сектора. Утверждение единственности имеет
форму
\begin{equation}
 \text{точечный центр и полная }SO(3)\text{-эквивариантность}
 \Longrightarrow q(\mathbf n)=[\mathbf n].
 \label{eq:v5-spin-bridge-conditional-chain}
\end{equation}
Само существование выделенного точечного центра и требование
эквивариантного граничного профиля ещё не выведены из вакуумного действия.
Инвариантность действия относительно вращений не означает, что каждое его
решение обязано быть точечно эквивариантным. Тривиальный однородный вакуум
может оставаться решением без проколотой сферы и без граничной карты.

Иными словами, круг устранён только наполовину:
\begin{enumerate}
 \item после появления сферически-симметричного точечного дефекта его
 хопфова линия и класс пятнадцать уже не назначаются вручную;
 \item почему вакуум обязан содержать хотя бы один такой точечный центр,
 текущая топология ещё не доказывает.
\end{enumerate}

\begin{nogocriterion}[Запрет кругового выбора граничного сектора]
Нельзя доказывать существование материи формулой $q(\mathbf n)=[\mathbf n]$,
если наличие проколотой сферы и нетривиального граничного условия уже
внесены в постановку. Единственность эквивариантного ежа доказывает форму
дефекта, но не самозарождение дефекта из однородного вакуума.
\end{nogocriterion}

\section{Статус и следующий тест}

Мост между ранними и поздними строгими блоками теперь имеет вид
\begin{equation}
 \boxed{
 S^2_{\rm sp}=SO(3)/SO(2)
 \xrightarrow{\;q\;}
 \mathbb{RP}^2_{\rm fam}=SO(3)/O(2)
 \xleftarrow{\;\pi\;}
 S^2_{\rm lift}}
\end{equation}
с подъёмами степени $\pm1$ и хопфовыми линиями $L/L^*$.

Следующий вопрос больше нельзя формулировать как поиск ещё одной линии или
индекса. Нужно проверить, выводит ли существующий родитель граничное
условие эквивариантного ежа или топологический суммарный закон, запрещающий
повсюду однородный сектор. Без такого закона глобальное принуждение материи
остаётся условным.

Следующий гейт:
\path{version5_equivariant_boundary_sector_selection_gate.tex}.

\begin{protocol}
\begin{enumerate}
 \item пространственная сфера направлений:
 \textbf{$SO(3)/SO(2)\simeq S^2$};
 \item пространство проекторных осей:
 \textbf{$SO(3)/O(2)\simeq\mathbb{RP}^2$};
 \item каноническая карта: \textbf{$q(n)=[n]$};
 \item $SO(3)$-эквивариантных карт этого типа: \textbf{одна};
 \item постоянная карта как полностью эквивариантный профиль:
 \textbf{нет};
 \item глобальных подъёмов выбранного ежа: \textbf{два};
 \item степени подъёмов: \textbf{$+1/-1$};
 \item классы хопфовых линий: \textbf{$+1/-1$};
 \item коэффициентные классы: \textbf{$+15/-15$};
 \item совместимость с $E/E^*$ и Real-обменом: \textbf{да};
 \item spin-cover bridge для эквивариантного ежа: \textbf{закрыт};
 \item вывод существования точечного центра из вакуума: \textbf{нет};
 \item безусловное исключение нулевого сектора: \textbf{ещё нет};
 \item следующий гейт: \textbf{выбор эквивариантного граничного сектора}.
\end{enumerate}
\end{protocol}

Машинный сертификат сохранён в
\path{s2t_v5_spin_cover_defect_sphere_bridge_gate_results.json}.